\documentclass[a4paper]{article}

\usepackage{amsfonts, amsmath, amssymb, amsthm}
\usepackage{fullpage}
\usepackage{hyperref}

\newtheorem{theorem}{Theorem}[section]
\newtheorem{lemma}[theorem]{Lemma}
\newtheorem{corollary}[theorem]{Corollary}
\newtheorem{proposition}[theorem]{Proposition}
\newtheorem{remark}[theorem]{Remark}
\newtheorem{example}[theorem]{Example}

\newcommand{\F}{\mathbb{F}}
\newcommand{\ord}[2][]{\operatorname{ord}_{#1}\!\left(#2\right)}
\newcommand{\abs}[1]{\left|#1\right|}
\renewcommand{\qedsymbol}{$\blacksquare$}

\setlength{\parindent}{0em}
\setlength{\parskip}{1em}

\begin{document}

\title{Finite-Support Stabilization for \textit{Lights Out} Nullity}
\author{William Boyles}
\date{\today}
\maketitle

\begin{abstract}
	Let $d(n)$ be the nullity over $\F_2$ of the $n\times n$ bounded \textit{Lights Out} grid.
	For every fixed finite set $S$ of odd primes, we give explicit thresholds $C_p(S)$ such that
	\[
		d\!\left(\prod_{p\in S}p^{e_p}-1\right)
		=
		d\!\left(\prod_{p\in S}p^{\min(e_p,C_p(S))}-1\right).
	\]
	Thus the nullity is eventually constant in each prime exponent.
	In particular, for every odd integer $a$, the sequence $d(a^k-1)$ is eventually constant.
	We also construct, for every $K\geq3$, an odd integer whose true smallest stabilization exponent is exactly $K$.
	The thresholds explain the observed transitions for $a=21$ and $a=57$, and reduce the repository's mixed-prime conjectures to small exact polynomial computations.
\end{abstract}

\section{Polynomial and order preliminaries}

Let $F_n(x)\in\F_2[x]$ be the Fibonacci polynomials
\[
	F_0(x)=0,\qquad F_1(x)=1,\qquad
	F_n(x)=xF_{n-1}(x)+F_{n-2}(x).
\]
Sutner's formula gives
\begin{equation}\label{eq:nullity}
	d(n)=\deg\gcd\!\left(F_{n+1}(x),F_{n+1}(x+1)\right)
\end{equation}
\cite{Sutner2000}.

For an odd integer $M>1$, define the exact-order stratum
\[
	R_M=\left\{\zeta+\zeta^{-1}:\ord{\zeta}=M\right\}.
\]
If $N$ is odd, the distinct roots of $F_N$ are the disjoint union
\begin{equation}\label{eq:strata}
	\bigcup_{\substack{M\mid N\\M>1}}R_M.
\end{equation}
Moreover, $F_N$ is the square of a square-free polynomial, so every root has multiplicity two \cite{HUNZIKER2004465}.

For odd $M$, let
\[
	\rho_2(M)=\min\{r>0:2^r\equiv\pm1\pmod M\}.
\]
If $\alpha\in R_M$, then
\begin{equation}\label{eq:degree}
	[\F_2(\alpha):\F_2]=\rho_2(M).
\end{equation}

Fix a finite set $S$ of odd primes.
For each $p\in S$, set
\begin{equation}\label{eq:thresholds}
	h_p=\ord[p]{2},
	\qquad
	c_p=v_p(2^{h_p}-1),
	\qquad
	a_p(S)=\max_{q\in S}v_p(h_q),
	\qquad
	C_p(S)=c_p+a_p(S).
\end{equation}
The term $c_p$ is the prime-power plateau length from the one-prime argument.
The additional term $a_p(S)$ measures how much $p$-power divisibility can enter through the orders belonging to the other primes in $S$.

For an odd integer $N$ supported on $S$, define its $S$-truncation by
\begin{equation}\label{eq:truncation}
	N^\flat=\prod_{p\in S}p^{\min(v_p(N),C_p(S))}.
\end{equation}

\section{Order growth after truncation}

\begin{lemma}\label{lem:order-truncation}
	Let $Q$ be an odd integer supported on $S$, let $Q^\flat$ be its truncation, and set
	\[
		m=\frac{Q}{Q^\flat}.
	\]
	Then
	\[
		\ord[Q]{2}=m\,\ord[Q^\flat]{2}.
	\]
\end{lemma}

\begin{proof}
	Write $u_p=v_p(Q)$.
	The Chinese remainder theorem and the prime-power order formula give
	\begin{equation}\label{eq:order-lcm}
		\ord[Q]{2}
		=
		\mathop{\operatorname{lcm}}_{\substack{p\in S\\u_p>0}}
		\left(h_p\,p^{\max(0,u_p-c_p)}\right).
	\end{equation}

	First consider a prime $\ell\notin S$.
	Its valuation in \eqref{eq:order-lcm} comes only from the fixed integers $h_p$.
	Truncation changes no prime support, so
	\[
		v_\ell\!\left(\ord[Q]{2}\right)
		=
		v_\ell\!\left(\ord[Q^\flat]{2}\right).
	\]

	Now fix $p\in S$.
	The $p$-adic valuation in \eqref{eq:order-lcm} is the maximum of
	\[
		\max(0,u_p-c_p)
		\qquad\text{and}\qquad
		\{v_p(h_q):q\in S,\ q\neq p,\ u_q>0\}.
	\]
	Every term in the second set is at most $a_p(S)$.
	If $u_p\leq C_p(S)$, truncation leaves $u_p$ unchanged, and truncating the other exponents does not change their fixed factors $h_q$.
	The $p$-adic valuation of the order is therefore unchanged.

	If $u_p>C_p(S)=c_p+a_p(S)$, then
	\[
		u_p-c_p>a_p(S),
	\]
	so the contribution from the $p$-power component strictly dominates all cross-prime contributions.
	Hence
	\[
		v_p\!\left(\ord[Q]{2}\right)=u_p-c_p.
	\]
	At the truncated exponent $C_p(S)$, the corresponding valuation is $a_p(S)$, and therefore
	\[
		v_p\!\left(\ord[Q]{2}\right)
		-
		v_p\!\left(\ord[Q^\flat]{2}\right)
		=
		u_p-C_p(S)
		=
		v_p(m).
	\]
	Comparing valuations at every prime proves the result.
\end{proof}

\section{Finite-support stabilization}

\begin{theorem}[Finite-support truncation]\label{thm:truncation}
	Let $S$ be a finite set of odd primes and let $N$ be an odd positive integer supported on $S$.
	Then
	\[
		d(N-1)=d(N^\flat-1),
	\]
	where $N^\flat$ is defined by \eqref{eq:truncation}.
\end{theorem}

\begin{proof}
	Suppose $\alpha$ and $\alpha+1$ are roots of $F_N$.
	By \eqref{eq:strata}, there are divisors $M,L>1$ of $N$ such that
	\[
		\alpha\in R_M,
		\qquad
		\alpha+1\in R_L.
	\]
	Because $\F_2(\alpha)=\F_2(\alpha+1)$, equation \eqref{eq:degree} gives
	\begin{equation}\label{eq:signed-equality}
		\rho_2(M)=\rho_2(L).
	\end{equation}

	Write $j_p=v_p(M)$ and $\ell_p=v_p(L)$.
	Fix $p\in S$ and suppose $j_p>C_p(S)$.
	As in the proof of Lemma \ref{lem:order-truncation}, the $p$-part belonging to $p^{j_p}$ dominates every cross-prime contribution, so
	\[
		v_p\!\left(\ord[M]{2}\right)=j_p-c_p>a_p(S).
	\]
	Signed order differs from ordinary order by at most a factor of two, so their valuations at the odd prime $p$ agree.
	Equation \eqref{eq:signed-equality} therefore forces
	\[
		v_p\!\left(\ord[L]{2}\right)=j_p-c_p.
	\]
	This valuation is larger than $a_p(S)$, so it cannot arise from any cross-prime order $h_q$.
	Consequently $\ell_p>C_p(S)$ and
	\[
		\ell_p-c_p=j_p-c_p,
		\qquad\text{hence}\qquad
		\ell_p=j_p.
	\]
	By symmetry, every exponent above its threshold occurs with the same value in $M$ and $L$.

	Set
	\[
		Q=\operatorname{lcm}(M,L),
		\qquad
		Q^\flat=\prod_{p\in S}p^{\min(v_p(Q),C_p(S))},
		\qquad
		m=\frac{Q}{Q^\flat}.
	\]
	If $m=1$, then $M,L\mid N^\flat$ and there is nothing to prove.
	Suppose $m>1$.

	Choose a primitive $Q$th root $\theta$.
	There are exponents $a,b$ such that
	\[
		\alpha=\theta^a+\theta^{-a},
		\qquad
		\alpha+1=\theta^b+\theta^{-b},
	\]
	where $\theta^a$ and $\theta^b$ have exact orders $M$ and $L$.
	Thus $\theta$ is a root of
	\begin{equation}\label{eq:sparse}
		P(x)=1+x^a+x^{Q-a}+x^b+x^{Q-b}.
	\end{equation}

	Every prime dividing $m$ occurs above its threshold in $Q$.
	The preceding rigidity argument says that its full exponent in $Q$ occurs in both $M$ and $L$.
	Since $\theta^a$ has exact order $M$, we have
	\[
		\gcd(Q,a)=\frac{Q}{M};
	\]
	no prime dividing $m$ divides $Q/M$, and no such prime can divide the remaining unit factor in $a$.
	Hence $\gcd(a,m)=1$.
	Likewise $\gcd(b,m)=1$.
	Because $m\mid Q$, only the constant exponent in \eqref{eq:sparse} is divisible by $m$.

	The element $\theta^m$ has exact order $Q^\flat$.
	Let $f(x)$ be its minimal polynomial over $\F_2$.
	Lemma \ref{lem:order-truncation} gives
	\[
		\deg f(x^m)
		=
		m\,\ord[Q^\flat]{2}
		=
		\ord[Q]{2},
	\]
	the degree of the minimal polynomial of $\theta$.
	Since $\theta$ is a root of $f(x^m)$, the polynomial $f(x^m)$ is its minimal polynomial.
	Therefore $f(x^m)\mid P(x)$.

	Decompose \eqref{eq:sparse} uniquely by exponents modulo $m$:
	\[
		P(x)=\sum_{r=0}^{m-1}x^rP_r(x^m).
	\]
	Divisibility by $f(x^m)$ implies that $f(y)$ divides every $P_r(y)$.
	But the preceding coprimality shows that $P_0(y)=1$, a contradiction.
	Thus $m=1$, and every common root of $F_N(x)$ and $F_N(x+1)$ already occurs for $F_{N^\flat}(x)$ and $F_{N^\flat}(x+1)$.

	The converse inclusion follows from $N^\flat\mid N$ and the Fibonacci-polynomial divisibility property.
	Because both indices are odd, every root has multiplicity two.
	Equation \eqref{eq:nullity} now gives
	\[
		d(N-1)=d(N^\flat-1).
	\]
\end{proof}

\begin{corollary}[Eventual constancy]\label{cor:eventual}
	For every odd positive integer $A$, there exists $k_0$ such that
	\[
		d(A^k-1)=d(A^{k_0}-1)
		\qquad\text{for every }k\geq k_0.
	\]
	More explicitly, write
	\[
		A=\prod_{p\in S}p^{b_p}
	\]
	and set
	\[
		k_0=
		\max_{p\in S}
		\left\lceil\frac{C_p(S)}{b_p}\right\rceil.
	\]
	Then this $k_0$ has the required property.
\end{corollary}

\begin{proof}
	For every $k\geq k_0$, the truncation of $A^k$ is the fixed integer
	\[
		B_S=\prod_{p\in S}p^{C_p(S)}.
	\]
	Theorem \ref{thm:truncation} gives
	\[
		d(A^k-1)=d(B_S-1).
	\]
	The same theorem applied to $A^{k_0}$ gives
	\[
		d(A^{k_0}-1)=d(B_S-1),
	\]
	which proves the claim.
\end{proof}

\section{Arbitrarily late stabilization}

The sufficient cutoff in Corollary \ref{cor:eventual} can be arbitrarily large.
More strongly, the actual smallest stabilization exponent is unbounded.

\begin{theorem}\label{thm:unbounded}
	For every integer $K\geq3$, there exists an odd positive integer $A_K$ such that $K$ is the smallest integer satisfying
	\[
		d(A_K^k-1)=d(A_K^K-1)
		\qquad\text{for every }k\geq K.
	\]
\end{theorem}

\begin{proof}
	Set
	\[
		r=3^{K-1},
		\qquad
		B=2^r+1.
	\]
	The lifting-the-exponent lemma gives
	\[
		v_3(B)=v_3(2+1)+v_3(r)=K.
	\]
	Write
	\[
		B=3^KQ,
		\qquad \gcd(Q,3)=1,
	\]
	and define
	\begin{equation}\label{eq:construction}
		A_K=3Q=\frac{2^{3^{K-1}}+1}{3^{K-1}}.
	\end{equation}

	Let $S$ be the prime support of $A_K$.
	If $q\mid Q$, then $2^r\equiv-1\pmod q$, so
	\[
		h_q=\ord[q]{2}\mid2r
		\qquad\text{but}\qquad
		h_q\nmid r.
	\]
	Since $r$ is a power of 3, this means
	\[
		h_q=2\cdot3^j
		\qquad\text{for some }0\leq j\leq K-1.
	\]
	Zsigmondy's theorem provides a primitive prime divisor of $2^{2r}-1$ because $2r\geq18$ \cite{Zsigmondy1892}.
	Such a prime has order exactly $2r$, so it divides $2^r+1$.
	It is not 3, whose order is 2, and therefore divides $Q$ and belongs to $S$.
	Hence the maximum value of $v_3(h_q)$ over $q\in S$ is $K-1$.
	Therefore
	\[
		C_3(S)=c_3+a_3(S)=1+(K-1)=K.
	\]

	Now let $q>3$ divide $Q$.
	Every $h_t$ for $t\in S$ divides $2r=2\cdot3^{K-1}$, so
	\[
		a_q(S)=0.
	\]
	Furthermore, $h_q\mid2r$ and $q\nmid 2r/h_q$.
	The lifting-the-exponent lemma gives
	\[
		c_q
		=
		v_q(2^{h_q}-1)
		=
		v_q(2^{2r}-1).
	\]
	Since $q\mid2^r+1$ and $q\nmid2^r-1$,
	\[
		c_q=v_q(2^r+1)=v_q(B).
	\]
	Consequently
	\[
		C_q(S)=v_q(B)
		\qquad(q\mid Q).
	\]

	The exponent of 3 in $A_K$ is one, while the exponent of every $q\mid Q$ is already $C_q(S)$.
	Hence Theorem \ref{thm:truncation} gives
	\begin{equation}\label{eq:constructed-sequence}
		d(A_K^k-1)
		=
		d(3^{\min(k,K)}Q-1).
	\end{equation}
	In particular the sequence is constant for $k\geq K$.

	It remains to show that it has not stabilized earlier.
	For $k<K$, set $N_k=3^kQ$.
	Then
	\[
		N_k=\frac{B}{3^{K-k}}\leq\frac{B}{3},
	\]
	and the elementary bound $d(n)\leq n$ gives
	\begin{equation}\label{eq:early-upper}
		d(A_K^k-1)=d(N_k-1)\leq N_k-1<\frac{B}{3}.
	\end{equation}

	On the other hand, the reduction and Kloosterman evaluations proved in
	\path{tex/upper_bound/upper_bound.tex} give
	\[
		d(2^r)=G(2^r+1)-2
	\]
	because $3\mid2^r+1$, and
	\[
		G(2^r+1)=\frac{2^r+1+\mathcal{K}_r}{2}.
	\]
	Consequently
	\[
		d(B-1)
		=
		d(2^r)
		=
		\frac{2^r-3+\mathcal{K}_r}{2},
	\]
	where $\mathcal{K}_r$ is the binary Kloosterman sum.
	Weil's bound $\abs{\mathcal{K}_r}\leq2\sqrt{2^r}$ \cite{Weil1948} yields
	\[
		d(B-1)
		\geq
		\frac{2^r-3-2\sqrt{2^r}}{2}.
	\]
	Because $r\geq9$,
	\[
		2^r-6\sqrt{2^r}-11>0,
	\]
	which is equivalent to
	\[
		\frac{2^r-3-2\sqrt{2^r}}{2}
		>
		\frac{2^r+1}{3}
		=
		\frac{B}{3}.
	\]
	Combining this with \eqref{eq:early-upper} shows
	\[
		d(A_K^k-1)<d(A_K^K-1)
		\qquad(k<K).
	\]
	Thus the true smallest stabilization exponent is exactly $K$.
\end{proof}

\begin{example}
	For $K=3$, construction \eqref{eq:construction} gives
	\[
		A_3=\frac{2^9+1}{3^2}=57,
	\]
	recovering the sequence
	\[
		0,\ 36,\ 252,\ 252,\ldots.
	\]
	For $K=4$ it gives
	\[
		A_4=\frac{2^{27}+1}{3^3}
		=4\,971\,027
		=3\cdot19\cdot87\,211,
	\]
	whose true stabilization exponent is 4.
\end{example}

\section{Examples}

\subsection{The family $3^n7^m$}

We have
\[
	h_3=2,\quad c_3=1,
	\qquad
	h_7=3,\quad c_7=1.
\]
Since $v_3(h_7)=1$, the thresholds for $S=\{3,7\}$ are
\[
	C_3(S)=2,
	\qquad
	C_7(S)=1.
\]
Therefore
\begin{equation}\label{eq:3-7-truncation}
	d(3^n7^m-1)
	=
	d\!\left(3^{\min(n,2)}7^{\min(m,1)}-1\right).
\end{equation}
The finite calculations
\[
	d(3\cdot7-1)=d(20)=0,
	\qquad
	d(3^2\cdot7-1)=d(62)=24
\]
now prove, for positive $m$,
\[
	d(3^n7^m-1)
	=
	\begin{cases}
		0 & n=1,\\
		24 & n\geq2.
	\end{cases}
\]
When $m=0$, prime-power stabilization gives $d(3^n-1)=0$.
Thus the second range in the repository conjecture should read $n\geq2$, not $n\geq1$.

For $A=21$, Corollary \ref{cor:eventual} gives $k_0=2$, and hence
\[
	d(21^k-1)=24
	\qquad(k\geq2).
\]

\subsection{The family $3^n19^m$}

Here
\[
	h_{19}=18,
	\qquad
	v_3(h_{19})=2,
\]
so
\[
	C_3(\{3,19\})=3,
	\qquad
	C_{19}(\{3,19\})=1.
\]
Consequently
\[
	d(3^n19^m-1)
	=
	d\!\left(3^{\min(n,3)}19^{\min(m,1)}-1\right).
\]
The exact finite values
\[
	d(57-1)=0,\qquad
	d(171-1)=36,\qquad
	d(513-1)=252
\]
prove, for positive $m$,
\[
	d(3^n19^m-1)
	=
	\begin{cases}
		0 & n=1,\\
		36 & n=2,\\
		252 & n\geq3.
	\end{cases}
\]
For $A=57$, the sufficient stabilization exponent is $k_0=3$:
\[
	d(57^k-1)=252
	\qquad(k\geq3).
\]

\subsection{The remaining two-prime conjectures}

The same thresholds reduce every positive-exponent two-prime formula listed in
\path{tex/conjectures/summary/conjectures.tex} to the following finite GCD calculations:
\begin{align*}
	d(15-1)&=4, &
	d(33-1)&=20, &
	d(39-1)&=d(117-1)=0,\\
	d(51-1)&=8, &
	d(35-1)&=4, &
	d(55-1)&=d(275-1)=4,\\
	d(65-1)&=28, &
	d(85-1)&=12, &
	d(95-1)&=4,\\
	d(77-1)&=0, &
	d(91-1)&=0, &
	d(119-1)&=8,\\
	d(133-1)&=0, &
	d(143-1)&=0, &
	d(187-1)&=8,\\
	d(209-1)&=0, &
	d(221-1)&=8, &
	d(247-1)&=0,\\
	d(323-1)&=8.
\end{align*}
These exact checks are performed by \texttt{code/multiprime\_stabilization\_check.py}.
Together with Theorem \ref{thm:truncation}, they prove all of those formulas for positive prime exponents.

\begin{remark}
	The thresholds $C_p(S)$ are sufficient but need not be minimal.
	For example, $C_5(\{5,11\})=2$ because $5\mid\ord[11]{2}$, but the two finite values at exponents one and two are both 4.
	Thus the nullity in that family is already constant from the first exponent.
\end{remark}

\bibliographystyle{amsalpha}
\bibliography{refs.bib}

\end{document}
